% Tracce uscite agli Esami di Stato (Sezione A) sugli argomenti del capitolo.
% Incluso da main.tex subito dopo il riassunto.

\section*{Tracce Uscite} % Usiamo * per non numerare la sezione nell'indice
\addcontentsline{toc}{section}{Tracce Uscite: Sistemi Operativi e Reti}

\subsection*{Seconda sessione 2018 -- Informatica}
\textbf{Traccia}: descrivere il modello ISO/OSI, esponendo nel dettaglio ogni livello e le sue caratteristiche; confrontarlo con il modello TCP/IP.

\textbf{Dove trovare la risposta}: capitolo \emph{Reti di Calcolatori} (sette strati) e, per il confronto con TCP/IP, sezione \emph{Lo Stack TCP/IP e i Principali Protocolli}, con la tabella di corrispondenza tra i livelli.

\subsection*{Prima sessione 2026 -- A2-INF-4}
\textbf{Traccia}: concetto di protocollo nei sistemi distribuiti; ruolo nel garantire interoperabilità, affidabilità e sincronizzazione; funzioni di un protocollo di comunicazione; modelli orientato e non orientato alla connessione con vantaggi, svantaggi e ambiti; esempi dallo stack TCP/IP e impiego nei sistemi distribuiti e nelle applicazioni Internet.

\textbf{Dove trovare la risposta}: sezioni \emph{Il Concetto di Protocollo}, \emph{Comunicazione Orientata e Non Orientata alla Connessione}, \emph{Lo Stack TCP/IP} e \emph{Protocolli nei Sistemi Distribuiti}, con la scaletta di risposta a fine capitolo.

\subsection*{Bologna (sessione non indicata nel file) -- Ing. Informatica E04}
\textbf{Traccia}: concetti di processo e thread in un sistema operativo moderno; differenze tra i due; modelli di multithreading (many-to-one, one-to-one, many-to-many); problematiche di concorrenza come le race condition.

\textbf{Dove trovare la risposta}: sezioni \emph{Processi: Stati, PCB e Cambio di Contesto}, \emph{Thread e Modelli di Multithreading} e \emph{Race Condition}; per le definizioni di base vedi anche il capitolo \emph{Sistemi Operativi}.

\subsection*{Prima sessione 2025 (Bologna) -- Ing. Informatica E05}
\textbf{Traccia}: processo di chiamata a funzione e ruolo dello stack.

\textbf{Dove trovare la risposta}: capitolo \emph{Programmazione ad Oggetti e Algoritmi}, sezione \emph{Chiamate di Funzione, Stack e Passaggio dei Parametri}.
